\documentclass[a4paper,12pt]{article}
\usepackage[utf8]{inputenc}
\usepackage{graphicx}
\graphicspath{ {./imgs/} }
\usepackage{fancyhdr} % For headers and footers
\usepackage{geometry}
\usepackage{hyperref} % For clickable links
\usepackage{longtable} % For long tables that can span multiple pages
\usepackage{array}     % For better column formatting

% Adjust margins
\geometry{
    top=4cm, 
    bottom=2.5cm, 
    left=2.5cm, 
    right=2.5cm,
    headheight=3cm, % Ensures sufficient space for header content
    headsep=0.5cm,
    }

% Define the custom headers
\fancypagestyle{firstpage}{% First page header style
    \fancyhf{} % Clear all header/footer fields
    \fancyhead[L]{%
        \includegraphics[height=1.5cm]{heig_logo.png} \\
    }
    \fancyhead[R]{%
        \includegraphics[height=1.5cm]{reds_logo.png} \\
    }
    \fancyfoot{} % No footer on the first page
}

% Define other pages style
\fancypagestyle{otherpages}{%
    \fancyhf{}
    \fancyhead[L]{%
        \includegraphics[height=1.5cm]{heig_logo.png} \\
    }
    \fancyhead[R]{%
        \small Laboratoire 7: Gestion de surcharge \\
        Rodrigo Lopez Dos Santos \\
        Urs Behrmann \\
    }
    \fancyfoot[C]{\small page \thepage}
}

\author{Urs Behrmann}

\setlength{\parskip}{1em} % Définit un espace entre chaque paragraphe
\setlength{\parindent}{0pt} % Optionnel : enlève l'indentation des paragraphes

\begin{document}

\emergencystretch=1em % Ajoute une flexibilité d'espacement

% First page content
\thispagestyle{firstpage}


\vspace*{2cm}
\begin{center}
    \Huge Laboratoire 7 \\
    \vspace{0.2cm}
    \Large Gestion de surcharge\\
    \vspace{1cm}
    \small Départements : TIC\\
    Unité d'enseignement PTR\\
\end{center}

\vspace{9cm}

\renewcommand{\arraystretch}{1.5} % Adjust row height

\begin{flushleft} % Left-align the table
    \begin{tabular}{@{}l l@{}}
        \textbf{Auteurs :}       & \textbf{Rodrigo Lopez Dos Santos} \\
                                 & \textbf{Urs Behrmann} \\
        \textbf{Professeur :}    & \textbf{Yorick Brunet} \\
        \textbf{Assistant :}     & \textbf{Anthony I. Jaccard} \\
        \textbf{Classe :}        & \textbf{PTR} \\
        \textbf{Salle de labo :} & \textbf{A09} \\
        \textbf{Date :}          & \textbf{28.05.2025} 
    \end{tabular}
\end{flushleft}

\newpage

% Apply the other pages style
\pagestyle{otherpages}

\tableofcontents

\newpage

\section{Introduction}



\newpage

\section{Etape 1 : Observation de la surcharge}

A une charge à 36, on commence à observer des 'overruns' dans le système. Ils ne sont pas encore très fréquents, mais ils sont présents.

\newpage

\section{Etape 2 : Détection de la surcharge avec watchdog}



\newpage

\section{Etape 3 : Fonctionnalité dégradée}



\newpage

\section{Conclusion}




\end{document}